\section*{Bab \ref{ch:g9:periodic-table-first-look} --- Tabel Periodik: Pandangan Pertama}

\begin{solution}{exo:g9:periodic-table-first-look:1}
Menurut \omterm{def:g9:inside-the-atom:atomic-number}{nomor atom} $Z$ yang makin besar.
\end{solution}

\begin{solution}{exo:g9:periodic-table-first-look:2}
Karbon: \omterm{def:g9:periodic-table-first-look:period}{periode} 2, \omterm{def:g9:periodic-table-first-look:period}{golongan} 14. Magnesium: \omterm{def:g9:periodic-table-first-look:period}{periode} 3, \omterm{def:g9:periodic-table-first-look:period}{golongan} 2.
Argon: \omterm{def:g9:periodic-table-first-look:period}{periode} 3, \omterm{def:g9:periodic-table-first-look:period}{golongan} 18.
\end{solution}

\begin{solution}{exo:g9:periodic-table-first-look:3}
\omterm{def:g9:periodic-table-first-look:metal}{Logam}: besi, tembaga, aluminium. \omterm{def:g9:periodic-table-first-look:metal}{Nonlogam}: belerang, oksigen, klorin.
\end{solution}

\begin{solution}{exo:g9:periodic-table-first-look:4}
Kalium (\ce{K}) di bawah natrium; bromin (\ce{Br}) di bawah klorin.
\end{solution}

\begin{solution}{exo:g9:periodic-table-first-look:5}
Kalsium (\ce{Ca}, $Z = 20$) dan nitrogen (\ce{N}, $Z = 7$).
\end{solution}

\begin{solution}{exo:g9:periodic-table-first-look:6}
Kalium berada dalam \omterm{def:g9:periodic-table-first-look:period}{golongan} yang sama dengan natrium, \omterm{def:g9:periodic-table-first-look:period}{golongan} 1, dan
unsur-unsur segolongan berperilaku sama.
\end{solution}

\begin{solution}{exo:g9:periodic-table-first-look:7}
\omterm{def:g9:periodic-table-first-look:metal}{Logam}: unsur itu mengilap, menghantarkan listrik, dan membentuk \omterm{def:g9:ions:ion}{ion}
positif.
\end{solution}

\begin{solution}{exo:g9:periodic-table-first-look:8}
Agar unsur-unsur yang sifatnya mirip tetap berada dalam kolom yang
sama, sekalipun harus membiarkan sebuah kotak kosong. Celah-celah itu
adalah ramalan: unsur-unsur yang hilang itu kemudian ditemukan dengan
sifat-sifat yang telah ia uraikan.
\end{solution}

\begin{solution}{exo:g9:periodic-table-first-look:9}
Karbon: $\ce{C} = 12$; oksigen: $\ce{O} = 16$. Celah ``? $= 68$'' (di
samping Al $= 27.4$) dan ``? $= 70$'' (di samping Si $= 28$).
\end{solution}

\begin{solution}{exo:g9:periodic-table-first-look:10}
\omterm{def:g9:periodic-table-first-look:period}{Periode} 2: 8 unsur (litium sampai neon). \omterm{def:g9:periodic-table-first-look:period}{Periode} 4: 18 unsur (kalium
sampai kripton).
\end{solution}

\begin{solution}{exo:g9:periodic-table-first-look:11}
\omterm{def:g9:periodic-table-first-look:period}{Golongan} 18. Di tempat yang tidak boleh ada reaksi apa pun: argon di
dalam bola lampu dan gas pengelasan, helium di dalam balon (helium tidak
\omterm{def:g4:what-a-fire-needs:burning}{terbakar}).
\end{solution}

\begin{solution}{exo:g9:periodic-table-first-look:12}
$(27.0 + 114.8)/2 = 70.9$, dekat dengan 69.7.
\end{solution}

\begin{solution}{pb:g9:periodic-table-first-look:1}
\textbf{1.} $Z = 32$; \omterm{def:g9:periodic-table-first-look:period}{periode} 4, \omterm{def:g9:periodic-table-first-look:period}{golongan} 14.

\textbf{2.} Di atas: silikon. Di bawah: timah. Di kiri: galium. Di
kanan: arsen.

\textbf{3.} Sebuah metaloid.

\textbf{4.} Germanium berada dalam \omterm{def:g9:periodic-table-first-look:period}{golongan} yang sama dengan silikon, dan
unsur-unsur segolongan berperilaku sama.

\textbf{5.} $(28.1 + 118.7)/2 = 73.4$.

\textbf{6.} $(69.7 + 74.9)/2 = 72.3$.

\textbf{7.} Rata-rata kiri--kanan (72.3): sepanjang \omterm{def:g9:periodic-table-first-look:period}{periode}, massa \omterm{def:g7:atoms-and-molecules:atom}{atom} relatif
naik secara teratur, unsur demi unsur, sedangkan ke bawah \omterm{def:g9:periodic-table-first-look:period}{golongan} massa
\omterm{def:g7:atoms-and-molecules:atom}{atom} melompat dengan langkah-langkah besar.

\textbf{8.} $72.6 - 70 = 2.6$.

\textbf{9.} $(28.1 + 118.7 + 69.7 + 74.9)/4 = 291.4/4 = 72.85$, yaitu
72.9.

\textbf{10.} Rata-rata itu, 72.9, dekat dengan angka Mendeleev, 72,
maupun angka Winkler, 72.3 (dan nilai sekarang, 72.6).

\textbf{11.} Sebuah unsur yang diramalkan sebelum ada yang pernah
melihatnya ditemukan dengan sifat-sifat yang telah diramalkan: tabel itu
bukan sekadar daftar, tabel itu dapat membuat ramalan.

\textbf{12.} Sekitar 72.9.
\end{solution}
